% Etapa 05 — ESQUELETO DE METODOLOGIA.
% Nao incluir no main.tex antes da revisao e integracao da Etapa 13.
% Planejamento: nenhuma execucao experimental foi realizada nesta etapa.
% Preservar o capitulo e a figura conceitual do TCC1.
\chapter{Metodologia e protocolo experimental}
\label{cap:metodologia-tcc2}

\section{Caracterizacao da investigacao}
\label{sec:met-caracterizacao}
% Descrever estudo experimental exploratorio de viabilidade tecnica.
% Objetivos aprovados em TCC1: NAO reescrever.
% Distinguir aplicacao Python PON de NOP Query C++.

\section{Prova de conceito e arquitetura investigada}
\label{sec:met-arquitetura}
% Descrever attribute-source, logic-engine e rule-executor,
% contratos HTTP/JSON, fluxo de notificacoes e gerenciamento de estado.
% Modelo 1: figura conceitual herdada do TCC1 em cap-proposta.tex.
% Modelo 2: diagrama especificado em docs/tcc2/etapa-04-arquitetura-distribuida.md.
% Inserir FIGURA DEFINITIVA do modelo implementado apos validacao
% nas Etapas 07--10, durante consolidacao da Etapa 13.
% Nao apresentar o modelo de projeto como implementacao comprovada.

\section{Ambiente experimental}
\label{sec:met-ambiente}
% PENDENTE Etapa 06: maquina, SO, CPU, RAM, virtualizacao,
% distribuicao Kubernetes, Docker, versoes, topologia e limites.
% Separar cluster virtual multinó de máquinas físicas.

\section{Cenarios e casos de teste}
\label{sec:met-cenarios}
% C0: processo unico; C1: tres processos locais;
% C2: Docker; C3: Kubernetes; C4: multinó condicional;
% C5: variacao de carga complementar; F1: reinicios condicionais.
% Matriz completa V01--V12 em docs/tcc2/etapa-05-protocolo-experimental.md.

\section{Entradas deterministicas e oraculo}
\label{sec:met-oraculo}
% Fixar fixture, estado inicial, seeds e hashes antes das medicoes.
% Comparar efeitos causais e ordem entre oraculo e cenarios distribuídos.

\section{Metricas e instrumentos de observacao}
\label{sec:met-metricas}
% Correcao funcional obrigatoria. Latencia, CPU, memoria, vazao e rede
% complementares quando houver instrumentacao valida.
% Relogios monotônicos so dentro da mesma origem.

\section{Procedimentos de coleta e reprodutibilidade}
\label{sec:met-reproducao}
% Descrever comandos efetivamente executados apenas apos experimentos,
% incluir SHA, digests, logs, manifestos, topologia, dados brutos e hashes.
% Etapa 11 executa; Etapa 12 analisa.

\section{Criterios de avaliacao e limitacoes}
\label{sec:met-criterios}
% V01--V08 e V11: evidencia funcional nuclear.
% V09 condicional; V10 recuperacao transparente; V12 complementar.
% Nao exigir melhoria de desempenho para atestar viabilidade.
